\includeonly{tex/research}
\documentclass{agh-wi}
\usepackage{scrhack}                        % Formatowanie kodów źródłowych programów

\titlePL{Zastosowanie fine-tuning'u dużych modeli językowych do wykrywania niechcianej poczty elektronicznej}  % Tytuł po polsku
\titleEN{Title consistent with the topic/field of the thesis} % Tytuł po angielsku
\author{Wojciech Dróżdż}
\supervisor{dr inż. Dariusz Żelasko}
 %%%%%%%%%%%%%%%%%%%%%%%%%%%%%%%%%%%%%
 % Pakiety
 %%%%%%%%%%%%%%%%%%%%%%%%%%%%%%%%%%%%%
 % Podstawowe — na pewno będziesz ich potrzebował(a) w swojej wersji pracy dyplomowej
 %%%%%%%%%%%%%%%%%%%%%%%%%%%%%%%%%%%%%
\usepackage[T1]{fontenc} % Naprawia problemy z niewspieranym encodinginem polskiego
\usepackage{lmodern} % Naprawia problemy z niewspieranym encodinginem polskiego
\usepackage{polski} % Obsługa języka polskiego
 %%%%%%%%%%%%%%%%%%%%%%%%%%%%%%%%%%%%%
 % Dodatkowe — niekoniecznie będziesz ich potrzebował(a) w swojej wersji pracy dyplomowej
 %%%%%%%%%%%%%%%%%%%%%%%%%%%%%%%%%%%%%
\usepackage[sorting=none]{biblatex}           % Spis literatury ma być tworzony na podstawie zawartości bibliograficznej bazy danych
\usepackage{amsmath}                          % Dodatkowe środowiska matematyczne
\usepackage{amssymb}                          % Dodatkowe symbole matematyczne
\usepackage[polish, intoc]{nomencl}           % Definiowanie symboli
\usepackage{graphicx}                         % Wstawianie grafik zewnętrznych
\usepackage{xcolor}                           % Kolorowy tekst
\usepackage{pgfplots}                         % Wykresy w LaTeX
\usepackage{xurl}                             % Lepsze łamanie długich adresów URL
\usepackage{tabularx}                         % Rozszerzona wersja środowiska 'tabular'
\usepackage{longtable}                        % Skład "długich" tabel
\usepackage[ruled,linesnumbered]{algorithm2e} % Algorytmy w formie pseudokodu
\usepackage{listings} 
\usepackage[newfloat]{minted}
% Sus jbc
\usepackage{float}
\makeatletter
\AddToHook{begindocument/before}{\@ifpackageloaded{minted}{\removefromtoclist[float]{lol}}{}}
\makeatother

\usepackage{hyperref}                         % Obsługa adresów URL — zbędny jeżeli używasz 'biblatex'
\usepackage{csquotes}                         % Polskie cudzysłowy — rekomendowany jeżeli  używasz 'biblatex'
\usepackage{framed}
\usepackage{xcolor}
\pgfplotsset{compat=1.18}
\usepgfplotslibrary{statistics}
\definecolor{plotHam}{HTML}{1B9E77}
\definecolor{plotSpam}{HTML}{D95F02}
\definecolor{plotBlue}{HTML}{3B528B}
\definecolor{plotAccent}{HTML}{7570B3}
\definecolor{plotGrid}{HTML}{D9DEE7}
 %%%%%%%%%%%%%%%%%%%%%%%%%%%%%%%%%%%%
 %  % Konfiguracja rysunków TikZ
 %%%%%%%%%%%%%%%%%%%%%%%%%%%%%%%%%%%%
\usetikzlibrary{shapes.geometric,arrows.meta} 
\tikzstyle{every picture}+=[remember picture]
\tikzstyle{na} = [baseline=-.5ex]
 %%%%%%%%%%%%%%%%%%%%%%%%%%%%%%%%%%%%
 % Ładowanie danych bibliograficznych
 %%%%%%%%%%%%%%%%%%%%%%%%%%%%%%%%%%%%
\addbibresource{bibliography.bib}
 %%%%%%%%%%%%%%%%%%%%%%%%%%%%%%%%%%%%%
 % Opcje konfiguracyjne pakietów
 %%%%%%%%%%%%%%%%%%%%%%%%%%%%%%%%%%%%%
 % Pakiet 'framed'
\definecolor{shadecolor}{gray}{0.9}
 % Pakiet 'nomencl'
\makenomenclature % Otwórz plik 'example.nlo'
 %%%%%%%%%%%%%%%%%%%%%%%%%%%%%%%%%%%%%
 % Definicje komend
 %%%%%%%%%%%%%%%%%%%%%%%%%%%%%%%%%%%%%
\newcommand{\alert}[1]{\colorbox{red!50}{#1}}
 %%%%%%%%%%%%%%%%%%%%%%%%%%%%%%%%%%%%%
 % Twierdzenia i podobne struktury
 %%%%%%%%%%%%%%%%%%%%%%%%%%%%%%%%%%%%%
\newtheorem{theorem}{Twierdzenie}
\newtheorem{definition}{Definicja}

\newcounter{comment}[chapter]
\newenvironment{comment}[1][]{\begin{shaded}\refstepcounter{comment}
   \noindent \textbf{Uwaga~\thechapter.\thecomment. #1} \rmfamily}{\end{shaded}}
\begin{document}
 %%%%%%%%%%%%%%%%%%%%%%%%%%%%%%%%%%%%%
\frontmatter % Część wstępna
 %%%%%%%%%%%%%%%%%%%%%%%%%%%%%%%%%%%%%
\maketitle % Dodaj stronę tytułową
 %%%%%%%%%%%%%%%%%%%%%%%
 % Jeżeli chcesz komuś podziękować, to możesz użyć poniższego kodu
\cleardoublepage
\thispagestyle{empty}
\vspace*{\fill}
\begin{flushright}
    \em
    \begin{minipage}{0.75\textwidth}
        Tutaj możesz umieścić treść podziękowań.
        Tutaj możesz umieścić treść podziękowań.
        Tutaj możesz umieścić treść podziękowań.
        Tutaj możesz umieścić treść podziękowań.
        Tutaj możesz umieścić treść podziękowań.
    \end{minipage}
\end{flushright}
 %%%%%%%%%%%%%%%%%%%%%%%
\begin{abstractPL}
    Streszczenie po polsku \ldots
\end{abstractPL}
\begin{abstractEN}
    Abstract in English \ldots
\end{abstractEN}
 %%%%%%%%%%%%%%%%%%%%%%%%%%%%%%%%%%%%%
 %%%%%%%%%%%%%%%%%%%%%%%%%%%%%%%%%%%%%
\tableofcontents   % Wygeneruj spis treści
\begin{shaded}
    Zawartość spisu treści \pauza tytuły rozdziałów oraz ich liczba zależą od tematyki pracy \pauza należy ustalić z opiekunem pracy.
\end{shaded}
\listoffigures     % Wygeneruj listę rysunków
\listoftables      % Wygeneruj listę tabel
\listofalgorithms  % Wygeneruj listę algorytmów
 % Wygeneruj listę kodów źródłowych
% \lstlistoflistings % Jeżeli do tworzenia listingów używasz pakietu 'listings'
 \listoflistings % Jeżeli do tworzenia listingów używasz pakietu 'minted'
\printnomenclature % Wyświetl listę symboli
 %%%%%%%%%%%%%%%%%%%%%%%%%%%%%%%%%%%%%
\mainmatter % Część główna
 %%%%%%%%%%%%%%%%%%%%%%%%%%%%%%%%%%%%%
 % Rozdziały części głównej
\chapter{Wstęp}

\section{Motywacja}

Poczta elektroniczna pozostaje jednym z podstawowych kanałów komunikacji
prywatnej i zawodowej. Jej powszechność pozwala niewielkim kosztem rozsyłać
wiadomości reklamowe, phishingowe i oszukańcze do dużej liczby odbiorców.
Filtr musi rozpoznawać zmieniające się sposoby formułowania wiadomości i nie
blokować poprawnej korespondencji. Fałszywy alarm może oznaczać przeoczenie
ważnej wiadomości, natomiast przepuszczenie ataku naraża odbiorcę na utratę
danych lub pieniędzy.

Jednym z etapów wielowarstwowego filtrowania poczty jest klasyfikacja tekstu.
W klasycznych metodach treść reprezentuje się między innymi za pomocą TF-IDF,
a następnie przekazuje do klasyfikatora. Modele językowe wykorzystują natomiast
reprezentacje zależne od kontekstu oraz wiedzę zdobytą podczas treningu
wstępnego. Ich zastosowanie wiąże się jednak z wyższym kosztem treningu i
inferencji niż w przypadku prostszych klasyfikatorów. Filtr pocztowy może
przetwarzać duży strumień wiadomości, dlatego oprócz skuteczności znaczenie mają
również czas odpowiedzi, przepustowość i zużycie zasobów. Powstaje zatem
pytanie, czy modele językowe uzyskują lepsze wyniki w klasyfikacji oraz jaki
jest ich koszt w porównaniu z klasycznymi klasyfikatorami.

\section{Cel pracy}

Celem pracy jest sprawdzenie, czy mały model językowy dostrojony za pomocą
adapterów LoRA może skutecznie klasyfikować wiadomości elektroniczne jako
\texttt{ham} albo \texttt{spam}. Badanie ma określić wpływ parametrów treningu
i sposobu sformułowania zadania na wyniki uzyskiwane na kilku korpusach
wiadomości.

Drugim celem jest porównanie dostrojonego modelu z modelem bez dostrajania oraz
prostszymi klasyfikatorami tekstu. Porównanie obejmuje skuteczność, profil
błędów i koszt inferencji.

\section{Zakres i metoda badawcza}

Badania dotyczą klasyfikacji wiadomości na podstawie tematu i treści. Nie
analizowano załączników, reputacji adresów i domen, pełnych nagłówków
transportowych ani wyników SPF, DKIM i DMARC. Przyjęta klasa \texttt{spam}
obejmuje niechciane reklamy, phishing oraz wiadomości oszukańcze. Model
rozstrzyga, czy wiadomość jest pożądana, ale nie określa rodzaju wykrytego
nadużycia.

Jako model bazowy wybrano \texttt{Qwen3-0.6B}, który dostrajano metodą LoRA.
Dla klasyfikacji przez ocenę następnego tokenu porównano 16 konfiguracji
różniących się długością kontekstu, współczynnikiem uczenia i parametrami
adaptera. Oceniono również pośrednie punkty kontrolne. Wybrane ustawienia
wykorzystano następnie do porównania czterech metod: klasyfikacji sekwencji,
generowania etykiety, oceny następnego tokenu oraz generowania odpowiedzi
ustrukturyzowanej.

Wyniki odniesiono do modelu bez dostrajania. Porównanie objęło Naive Bayes,
regresję logistyczną i SVM korzystające z reprezentacji TF-IDF, a także
\texttt{fastText}, MiniLM z regresją logistyczną oraz DistilBERT. Dla metod
bazowych dobrano parametry i progi klasyfikacji na wydzielonym zbiorze
walidacyjnym. Końcowe porównanie przeprowadzono na czterech zbiorach testowych,
po 2000 wiadomości w każdym. Koszt inferencji zmierzono na wspólnej próbce 1000
wiadomości.

\section{Pytania badawcze}

Zakres pracy opisują następujące pytania badawcze:

\begin{enumerate}
    \item Jakie wyniki osiąga mały model językowy dostrojony za pomocą LoRA na
    próbce treningowej, w zewnętrznej walidacji i w teście końcowym?
    \item Który z badanych sposobów użycia modelu językowego daje
    najkorzystniejsze połączenie skuteczności, niezależności decyzji od formatu
    generowanej odpowiedzi i prostoty inferencji?
    \item Jak parametry treningu oraz jego czas trwania wpływają na wynik
    modelu na korpusach niewykorzystanych w treningu?
    \item Jaki kompromis między skutecznością a kosztem inferencji wynika z
    porównania dostrojonego modelu językowego z prostszymi klasyfikatorami?
\end{enumerate}

\section{Struktura pracy}

Po wstępie przedstawiono podstawy teoretyczne klasyfikacji wiadomości oraz
działania modeli językowych. Rozdział obejmuje tokenizację, architekturę
Transformer, dostrajanie adapterowe i sposoby wykorzystania modelu językowego
do klasyfikacji. Następnie omówiono mechanizmy stosowane w filtrowaniu poczty,
od reguł i reputacji po klasyczne algorytmy uczenia maszynowego oraz
klasyfikatory transformerowe.

Część badawcza opisuje wybór modelu, przygotowanie zbiorów danych i
przebieg eksperymentów. Kolejne sekcje dotyczą doboru parametrów LoRA, analizy
punktów kontrolnych, porównania sposobów dostrajania oraz zestawienia
wariantów Qwen3 z metodami bazowymi i modelem bez dostrajania na wspólnym
teście końcowym. Rozdział końcowy odpowiada na pytania badawcze, omawia
ograniczenia, możliwość wdrożenia klasyfikatora i dalsze kierunki badań.

\chapter{Przegląd metod wykrywania niechcianej poczty}

Praktyczny filtr pocztowy korzysta nie tylko z tematu i treści wiadomości, ale
też z nagłówków SMTP, adresów IP, domen, historii nadawcy, wyników
uwierzytelniania i informacji o odsyłaczach. Analiza tekstu jest jedną z jego
warstw. Pozostałe mechanizmy oceniają nadawcę, wykrywają znane kampanie albo
sprawdzają podejrzane adresy URL. W przeglądzie uwzględniono zarówno te
rozwiązania, jak i metody możliwe do uruchomienia na zbiorach tekstowych użytych
w części badawczej.

\section{Wielowarstwowe systemy filtrowania poczty}

W produkcyjnych skrzynkach pocztowych pierwsza decyzja często zapada jeszcze przed
analizą pełnej treści wiadomości. Dostawca poczty może sprawdzić, czy nadawca
spełnia wymagania techniczne, czy wiadomość przechodzi uwierzytelnianie SPF,
DKIM i DMARC, czy adres IP ma dobrą reputację oraz czy domena nie występuje w
kampaniach nadużyć. Wymagania publikowane dla nadawców poczty Gmail wskazują
wprost na znaczenie SPF lub DKIM, polityki DMARC, poprawnego formatu wiadomości
i utrzymywania niskiego odsetka zgłoszeń spamu~\autocite{gmail-sender-guidelines}.
Podobny obraz widać w ustawieniach Microsoft Defender for Office 365, gdzie
osobno konfigurowane są między innymi polityki antyspamowe, antyphishingowe,
obsługa podszywania się pod nadawcę i reakcje na DMARC~\autocite{microsoft-defender-o365-security-settings}.

SpamAssassin opiera decyzję na zestawie testów punktowanych, regułach
treści, listach dozwolonych i blokowanych, testach sieciowych oraz mechanizmach
uczących się~\autocite{spamassassin-conf}. Rspamd udostępnia osobne moduły dla
SPF, DKIM, DMARC, reputacji, list DNS, fuzzy hashy, phishingu, sieci neuronowych
i wielu innych sygnałów~\autocite{rspamd-modules}. Oba filtry działają jako
systemy punktowe, które składają wiele sygnałów w jedną decyzję.

\section{Metody regułowe, reputacyjne i uwierzytelnianie nadawcy}

W filtrze pocztowym reguła jest jawnie zapisanym warunkiem sprawdzanym na
wiadomości lub jej metadanych. Po spełnieniu warunku filtr może dodać określoną
liczbę punktów do wyniku spamu, zablokować wiadomość albo przekazać ją do dalszej
analizy. Reguły są nadal używane, ponieważ działają szybko, są interpretowalne i
pozwalają reagować na znane schematy nadużyć bez ponownego trenowania modelu.
Reguła może dotyczyć fragmentu nagłówka, nietypowego adresu URL, słowa często
występującego w danej kampanii albo kombinacji kilku prostszych testów.
Administrator może też podnieść lub obniżyć wagę konkretnego testu, co ułatwia
bezpośrednią kontrolę nad polityką filtrowania.

Słabą stroną reguł jest zależność od wcześniej zauważonych wzorców. Nadawcy
spamu mogą zmieniać pisownię słów, przenosić treść do obrazów, rotować domeny
lub modyfikować szablony wiadomości. Dlatego reguły zwykle nie występują same,
lecz są łączone z reputacją nadawcy, listami domen, uwierzytelnianiem i
klasyfikatorami treści. Przeglądy metod antyphishingowych również pokazują, że
obrona przed phishingiem obejmuje kilka rodzin podejść: listy znanych adresów,
heurystyki, analizę treści, analizę wizualną oraz metody uczenia maszynowego
~\autocite{wood-antiphishing-review-2022}.

Uwierzytelnianie nadawcy rozwiązuje inny problem niż klasyfikacja tekstu.
Mechanizmy SPF, DKIM i DMARC pomagają ustalić, czy wiadomość mogła zostać
wysłana w imieniu deklarowanej domeny. Poprawnie uwierzytelniona wiadomość może
być niechcianą reklamą albo pochodzić z przejętego konta, dlatego analiza treści
pozostaje potrzebna.

\section{Klasyczne metody uczenia maszynowego}

Klasyczne modele uczenia maszynowego oparte na cechach leksykalnych są
najprostszym porównaniem dla klasyfikacji treści. Wiadomość zamienia się wtedy
na wektor cech, na przykład przez zliczanie słów, n-gramy albo TF-IDF, a
następnie trenuje klasyfikator taki jak Naive Bayes, regresja logistyczna lub
SVM. Jest to stosunkowo proste podejście, które zapewnia tani trening i szybką
inferencję, a przy odpowiednio dobranych danych może osiągać dobre wyniki.

W pracy \emph{Learning to Filter Spam E-Mail: A Comparison of a Naive Bayesian
and a Memory-Based Approach} porównano Naive Bayes z podejściem pamięciowym.
Automatycznie uczone filtry wyraźnie przewyższały ręcznie konstruowane reguły
słów kluczowych~\autocite{androutsopoulos-naive-bayes-spam}. Praca
\emph{Classification of Spam Emails through Hierarchical Clustering and
Supervised Learning} obejmowała porównanie reprezentacji TF-IDF i bag-of-words
z klasyfikatorami Naive Bayes, drzewami decyzyjnymi oraz SVM
~\autocite{janez-martino-spam-supervised-2020}.

Porównanie z klasycznymi metodami pozwala ocenić, ile z wyniku można uzyskać
dzięki dopasowaniu do słownictwa zbioru treningowego. Wynik TF-IDF bliski
wynikowi modelu językowego podważałby zasadność użycia droższego rozwiązania.
Przewaga modelu językowego na danych z innych źródeł wskazywałaby natomiast na
lepszą generalizację.

Reprezentacja leksykalna dobrze wychwytuje powtarzalne zwroty, nazwy produktów,
typowe frazy oszustw i często występujące domeny. Gorzej radzi sobie z
parafrazą, długim kontekstem i wiadomościami, w których znaczenie wynika z
relacji między kilkoma zdaniami, ponieważ warianty typu bag-of-words tracą
kolejność słów i słabiej opisują zależności semantyczne
~\autocite{le-distributed-representations-2014}. Rzadkie cechy oraz krótka treść
wiadomości mogą dodatkowo prowadzić do przeuczenia i słabszej generalizacji
~\autocite{xu-alternative-text-representation-2013}. W danych e-mailowych oznacza
to ryzyko dopasowania się do cech konkretnego korpusu, na przykład
charakterystycznego formatowania albo powtarzalnych fragmentów stopki.

\section{Szybkie modele tekstowe i reprezentacje semantyczne}

Pomiędzy klasycznym TF-IDF a pełnym dostrajaniem transformera znajdują się metody
pośrednie. Jedną z nich jest \texttt{fastText}, który uczy reprezentacje słów z
n-gramów znakowych i stosuje liniowy klasyfikator. Reprezentacje słów
występujących w dokumencie są łączone przed wykonaniem klasyfikacji. Autorzy
\texttt{fastText} pokazali, że taki model może być konkurencyjny wobec
głębszych klasyfikatorów tekstu, a jednocześnie trenuje się i działa bardzo
szybko~\autocite{joulin-fasttext-bag-of-tricks}.
W eksperymentach użyto klasyfikatora nadzorowanego z n-gramami znakowymi
długości 3--6, unigramami słów, wektorami o wymiarze 100, współczynnikiem
uczenia 0,25 i 20 epokami treningu.

Drugim wariantem pośrednim jest użycie gotowego modelu embeddingowego do
zamiany wiadomości na wektor semantyczny, a następnie wytrenowanie prostego
klasyfikatora liniowego. Sentence-BERT został zaproponowany właśnie po to, aby
uzyskiwać reprezentacje zdań nadające się do porównań semantycznych i dalszych
zadań bez kosztownego porównywania każdej pary tekstów przez pełny model BERT
~\autocite{reimers-sentence-bert}. W eksperymentach sprawdzono ten wariant za
pomocą MiniLM i regresji logistycznej.

Słabszym punktem takiego wariantu jest brak treningu pod konkretny typ nadużyć.
Gotowy model embeddingowy może dobrze oddawać ogólne znaczenie tekstu, ale
słabiej reagować na drobne sygnały istotne dla filtra pocztowego: nietypowe
domeny, zapis kwot, celowe literówki albo powtarzalne elementy szablonów
kampanii.

\section{Modele transformerowe i językowe}

Modele transformerowe budują reprezentację zależną od kontekstu. W wariancie
klasyfikacyjnym model, taki jak BERT lub DistilBERT, otrzymuje tekst wiadomości
i zwraca etykietę przez głowicę klasyfikacyjną. DistilBERT jest mniejszą wersją
BERT-a uzyskaną przez destylację wiedzy; według autorów zachowuje dużą część
jakości modelu bazowego, przy mniejszym rozmiarze i szybszym działaniu
~\autocite{sanh-distilbert}.

Modele językowe wymagają więcej pamięci niż metody TF-IDF, a ich inferencja jest
wolniejsza i zależy od długości kontekstu. Mogą jednak lepiej wykorzystywać
kontekst i uogólniać poza dosłowne słownictwo zbioru treningowego. Dlatego w
części badawczej oceniono również ich profil błędów na danych spoza zbioru
treningowego.

\subsection{Badania nad wykorzystaniem modeli językowych w filtrowaniu poczty}

W badaniach nad użyciem modeli językowych do filtrowania poczty stosowano
dostrajanie na oznaczonych wiadomościach, klasyfikację na podstawie promptu oraz
analizę treści uzupełnionej o nagłówki lub reputację adresów URL.

W pracy \emph{Spam-T5: Benchmarking Large Language Models for Few-Shot Email
Spam Detection} porównano klasyczne klasyfikatory, RoBERTę, SetFit oraz Spam-T5,
czyli Flan-T5 trenowany do generowania etykiety wiadomości. Eksperyment
obejmował cztery korpusy i próby treningowe od 4 przykładów do pełnego zbioru.
Spam-T5 uzyskał najlepsze wyniki przy 4--16 przykładach oraz najwyższy średni
F1 liczony dla wszystkich badanych rozmiarów próby. Trening i inferencja trwały
dłużej niż w przypadku klasycznych metod~\autocite{labonne-spam-t5-2023}.
W pracy \emph{Next-Generation Spam Filtering: Comparative Fine-Tuning of LLMs,
NLPs, and CNN Models for Email Spam Classification} dostrojono GPT-4, BERT-a i
RoBERTę na dwóch zbiorach wiadomości. Po losowym podziale danych modele osiągały
accuracy przekraczające 98\%. Ocena na drugim korpusie przyniosła niższe wyniki,
szczególnie dla BERT-a i RoBERT-y
~\autocite{roumeliotis-next-generation-spam-2024}.

Klasyfikacja przez prompt wykorzystuje gotowy model wraz z instrukcją i
przykładami umieszczonymi w kontekście. W pracy \emph{Evaluating the Performance
of ChatGPT for Spam Email Detection} sprawdzono ChatGPT w trybie zero-shot,
one-shot i five-shot. Na większym zbiorze anglojęzycznym najlepszy wariant
ChatGPT uzyskał macro F1 równe 80\%, a nadzorowany BERT 98\%. Na małym zbiorze
chińskojęzycznym różnica między ChatGPT a metodami nadzorowanymi była mniejsza
~\autocite{si-chatgpt-spam-2024}. W pracy \emph{Zero-Shot Spam Email
Classification Using Pre-trained Large Language Models} oceniono ChatGPT, GPT-4
i Flan-T5 na całym korpusie SpamAssassin. Modele klasyfikowały skróconą treść
albo streszczenie utworzone wcześniej przez ChatGPT. Etap streszczania podniósł
F1 GPT-4 z 88\% do 95\% i wymagał dodatkowego wywołania modelu. Badanie objęło
jeden korpus, SpamAssassin~\autocite{rojas-galeano-zero-shot-spam-2024}.

W pracy \emph{ChatSpamDetector: Leveraging Large Language Models for Effective
Phishing Email Detection} opisano system przetwarzający nagłówki i elementy
HTML, który oprócz klasy zwracał wyjaśnienie decyzji. GPT-4 osiągnął accuracy
równe 99,70\%. Średni czas jednego zapytania wyniósł 12,53~s, a ocena 2010
wiadomości kosztowała 265,80~USD według cen API użytych w badaniu
~\autocite{koide-chatspamdetector-2024}. W pracy \emph{APOLLO: A GPT-Based Tool
to Detect Phishing Emails and Generate Explanations That Warn Users} połączono
GPT-4o z reputacją adresów URL. Poprawne dane zewnętrzne podniosły accuracy z
97,4\% do 99,9\%. Po podaniu błędnych informacji o adresach URL wynik spadł
poniżej 46\%~\autocite{desolda-apollo-2025}.

Wartości zestawione w tabeli~\ref{tab:related-llm-email-detection} pochodzą z
różnych zbiorów, podziałów danych i zakresów informacji wejściowych. Ich
bezpośrednie porównanie ma charakter orientacyjny.

\begin{table}[!htbp]
\centering
\scriptsize
\resizebox{\textwidth}{!}{\begin{tabular}{p{0.18\textwidth}p{0.21\textwidth}p{0.23\textwidth}p{0.31\textwidth}}
\hline
\textbf{Praca} & \textbf{Modele i sposób użycia} & \textbf{Dane} & \textbf{Najważniejszy wynik i ograniczenie} \\
\hline
Labonne i Moran~\autocite{labonne-spam-t5-2023} & Spam-T5, RoBERTa, SetFit i metody klasyczne; trening pełny i few-shot & Ling-Spam, SMS Spam Collection, SpamAssassin i Enron & F1 Spam-T5 95--99\% przy pełnym treningu; przewaga przy 4--16 przykładach, lecz wyższy czas treningu i inferencji \\
Si i in.~\autocite{si-chatgpt-spam-2024} & ChatGPT zero-, one- i five-shot; BERT oraz metody klasyczne & angielski ESD i mały zbiór chińskojęzyczny & macro F1 80\% dla ChatGPT i 98\% dla BERT-a na ESD; odwrotna relacja na małym zbiorze chińskim \\
Rojas-Galeano~\autocite{rojas-galeano-zero-shot-spam-2024} & ChatGPT, GPT-4 i Flan-T5 bez dostrajania & SpamAssassin, 6047 wiadomości & F1 90\% dla Flan-T5 na treści i 95\% dla GPT-4 na streszczeniach; jeden korpus i dodatkowy etap generowania \\
Koide i in.~\autocite{koide-chatspamdetector-2024} & GPT-4, GPT-3.5, Llama~2 i Gemini Pro; treść oraz nagłówki & 1010 wiadomości phishingowych i 1000 pożądanych & accuracy 99,70\% dla GPT-4; średnio 12,53~s na wiadomość i 265,80~USD za ocenę zbioru \\
Desolda i in.~\autocite{desolda-apollo-2025} & GPT-4o, treść wiadomości i dane o adresach URL & 2000 wiadomości phishingowych i 2000 pożądanych & accuracy 97,4\% bez danych zewnętrznych i 99,9\% z poprawną reputacją URL; silna zależność od jakości tych danych \\
Roumeliotis i in.~\autocite{roumeliotis-next-generation-spam-2024} & dostrajanie GPT-4, BERT-a, RoBERT-y i CNN & dwa zbiory z serwisu Kaggle & accuracy powyżej 98\% na własnych testach; spadek przy ocenie na drugim zbiorze \\
\hline
\end{tabular}
}
\caption{Badania nad wykorzystaniem modeli transformerowych i językowych do
klasyfikacji spamu oraz phishingu w poczcie elektronicznej.}
\label{tab:related-llm-email-detection}
\end{table}
\FloatBarrier

W analizowanych pracach małe modele generatywne z otwartymi wagami, liczące
poniżej miliarda parametrów, pojawiają się rzadko. Badania skupiają się zwykle
na jednym sposobie użycia modelu. Spośród zestawionych prac tylko
\emph{Next-Generation Spam Filtering} obejmuje test przenoszenia modelu między
korpusami.

Wkład własny pracy obejmuje porównanie czterech sposobów realizacji
klasyfikacji przez Qwen3-0.6B przy zachowaniu wspólnego modelu, danych i
protokołu oceny. Porównanie przeprowadzono na kilku korpusach, a wyniki
zestawiono z prostszymi klasyfikatorami oraz kosztem inferencji wszystkich
metod.

\section{Zakres porównania eksperymentalnego}

Bezpośrednie porównanie objęło metody korzystające z treści wiadomości.
Wiadomości przypisywano do klasy \texttt{ham} albo \texttt{spam}. Wykorzystane
zbiory nie zawierają spójnych danych o reputacji nadawcy, DNS, pełnych
nagłówkach ani historii ruchu.
Tabela~\ref{tab:sota-method-scope} wskazuje, które metody omówiono jako kontekst
literaturowy, a które uruchomiono w eksperymentach.

\begin{table}[H]
\centering
\small
\resizebox{\textwidth}{!}{\begin{tabular}{p{0.22\textwidth}p{0.31\textwidth}p{0.18\textwidth}p{0.18\textwidth}}
\hline
\textbf{Grupa metod} & \textbf{Główne źródło sygnału} & \textbf{Wymaga metadanych poczty} & \textbf{Ujęcie w eksperymentach} \\
\hline
Uwierzytelnianie i reputacja nadawcy & Nagłówki SMTP, DNS, SPF, DKIM, DMARC, reputacja adresów IP i domen & Tak & Nie \\
Systemy regułowe i reputacyjne & Reguły treści, adresy URL, listy reputacyjne, scoring, czasem filtr bayesowski & Częściowo & Nie \\
TF-IDF + klasyfikator & Treść wiadomości reprezentowana przez cechy leksykalne & Nie & Tak \\
\texttt{fastText} & N-gramy słów i szybki model liniowy z embeddingami & Nie & Tak \\
Embeddingi zdań + klasyfikator & Wektor semantyczny wiadomości i klasyfikator liniowy & Nie & Tak \\
Transformer klasyfikacyjny & Dostrajany enkoder tekstu z głowicą klasyfikacyjną & Nie & Tak \\
Mały model językowy z LoRA & Dostrojony model generatywny, etykieta lub prawdopodobieństwo etykiety & Nie & Główna metoda opisywana w pracy \\
\hline
\end{tabular}
}
\caption{Zakres metod omawianych w przeglądzie oraz ich wykorzystanie w części
eksperymentalnej.}
\label{tab:sota-method-scope}
\end{table}

\chapter{Część badawcza}
    \section{Wybór otwartoźródłowego modelu językowego}
        \subsection{Przegląd dostępnych modeli (np. Llama, Mistral, Falcon, itp.)}
        \subsection{Kryteria wyboru modelu (wydajność, rozmiar, licencja, możliwości fine-tuning'u)}
        \subsection{Uzasadnienie wyboru konkretnego modelu}
    \section{Metodologia badawcza}
        \subsection{Opis wykorzystanych zbiorów danych}
            \subsubsection{Zbiór treningowy}
            \subsubsection{Zbiór walidacyjny}
            \subsubsection{Zbiór testowy}
        \subsection{Przygotowanie danych (Preprocessing)}
        \subsection{Proces fine-tuning'u modelu}
            \subsubsection{Architektura zadania klasyfikacyjnego}
            \subsubsection{Parametry i hiperparametry uczenia}
            \subsubsection{Wykorzystane narzędzia i biblioteki}
        \subsection{Metryki oceny skuteczności}
            \subsubsection{Dokładność (Accuracy)}
            \subsubsection{Precyzja (Precision)}
            \subsubsection{Pełność (Recall)}
            \subsubsection{Miara F1 (F1-Score)}
            \subsubsection{Inne istotne metryki (np. AUC)}
    \section{Implementacja i przeprowadzenie eksperymentów}
        \subsection{Środowisko eksperymentalne}
        Witajcie co tam  HALo!
        \subsection{Przebieg procesu fine-tuning'u}
        KUUURWA
    \section{Analiza skuteczności dostrojonego modelu}
        \subsection{Wyniki na zbiorze testowym}
        \subsection{Porównanie z wynikami istniejących metod (z Rozdziału 2)}
        \subsection{Analiza błędów klasyfikacji (False Positives, False Negatives)}
        \subsection{Ocena elastyczności modelu i odporności na nowe typy spamu} % Odporność na ewolucję zagrożeń
    \section{Analiza możliwości wdrożenia produkcyjnego}
        \subsection{Wymagania sprzętowe i czas wnioskowania (inference time)}
        \subsection{Koszty obliczeniowe (fine-tuning i wnioskowanie)}
        \subsection{Skalowalność rozwiązania}
        \subsection{Potencjalne wyzwania i ograniczenia wdrożeniowe}
        \subsection{Praktyczność zastosowania w porównaniu do tradycyjnych filtrów}

\chapter{Podsumowanie i wnioski}

\section{Synteza przeprowadzonych badań}

Praca dotyczyła wykorzystania modelu Qwen3-0.6B do wykrywania niechcianej
poczty na podstawie tematu i treści wiadomości. Publiczne korpusy sprowadzono
do wspólnego schematu, usunięto puste rekordy i duplikaty, a dane treningowe
zbalansowano. Trzy źródła pozostawiono poza treningiem, aby oceniać model także
na wiadomościach o innej charakterystyce.

Model bez dostrajania przypisywał wszystkie oceniane wiadomości do klasy
\texttt{ham}. Na zbalansowanym \texttt{train\_subset} oznaczało to 50,0\%
\emph{accuracy} i F1 równe 0 dla klasy pozytywnej. Po dostrojeniu najlepszy
wynik uzyskała konfiguracja \texttt{K08}: kontekst 512 tokenów, współczynnik
uczenia \(10^{-4}\), LoRA \(r=32\), \(\alpha=64\) oraz dropout 0,05. Siódmy
oceniany punkt kontrolny osiągnął \(S_{\mathrm{ext}}=98{,}0\%\), F1 równe
97,2\% na \texttt{spam\_ham}, \emph{specificity} 96,9\% na \texttt{enron} i
\emph{recall} 99,9\% na \texttt{fraudulent\_email\_corpus}.

Najlepsze warianty generowania etykiety, oceny następnego tokenu i odpowiedzi
ustrukturyzowanej różniły się wynikiem zbiorczym o około 0,1 punktu
procentowego. Najlepszą metodą bazową był DistilBERT z wynikiem
\(S_{\mathrm{ext}}=92{,}3\%\). TF-IDF z Naive Bayes i SVM osiągnęły
odpowiednio 90,7\% i 90,1\%.

\section{Odpowiedzi na pytania badawcze}

\begin{enumerate}
    \item \textbf{Czy dostrojenie małego modelu językowego za pomocą LoRA
    pozwala poprawnie klasyfikować wiadomości ze zbioru treningowego i ze źródeł
    niewykorzystanych do aktualizacji wag?}

    Sam prompt nie wystarczył, aby model bez dostrajania poprawnie wykonywał
    klasyfikację. Po dostrojeniu z wykorzystaniem adaptera LoRA model uzyskał
    wysokie wyniki na podzbiorze danych treningowych oraz na trzech innych
    korpusach.

    \item \textbf{Który z badanych sposobów użycia modelu językowego daje
    najkorzystniejsze połączenie skuteczności, niezależności decyzji od formatu
    generowanej odpowiedzi i prostoty inferencji?}

    Do dalszego porównania wybrano metodę oceny następnego tokenu. Jej przewaga
    liczbowa nad generowaniem etykiety i odpowiedzi ustrukturyzowanej była
    niewielka. Za wyborem przemawiał brak generowania tekstu: decyzja wynikała
    bezpośrednio z porównania prawdopodobieństw etykiet i nie zależała od
    parsera. Klasyfikacja sekwencji również nie wymagała parsera, ale uzyskała
    niższy wynik zewnętrzny, równy 96,5\%.

    \item \textbf{Jak parametry treningu oraz jego czas trwania wpływają
    na wynik modelu na korpusach niewykorzystanych w treningu?}

    Metryki na danych podobnych do treningowych szybko zbliżały się do 100\%,
    przez co słabo rozróżniały konfiguracje. Warianty ze współczynnikiem uczenia
    \(10^{-6}\) uzyskały niższe wyniki na pozostałych korpusach, a kontekst 1024
    tokenów wydłużał trening bez poprawy najlepszego rezultatu. W grupie
    wariantów ze współczynnikiem \(10^{-4}\) najlepiej wypadł adapter z
    \(r=32\) i \(\alpha=64\). Mniejszy adapter uzyskał słabszy wynik, natomiast
    zwiększenie pojemności ponad \(r=32\) i \(\alpha=64\) nie dawało regularnej
    poprawy. Nie udało się określić jednoznacznego wpływu dropout. Dla
    najlepszych konfiguracji najwyższe \(S_{\mathrm{ext}}\) pojawiało się przed
    końcem treningu. Wybór
    ostatniego adaptera lub kierowanie się wyłącznie metrykami treningowymi
    prowadziłoby do wyboru słabszego wariantu.

    \item \textbf{Jaki kompromis między skutecznością a kosztem inferencji
    wynika z porównania dostrojonego modelu językowego z prostszymi
    klasyfikatorami?}

    Qwen3 uzyskał wynik \(S_{\mathrm{ext}}\) wyższy o 5,7 punktu procentowego
    od DistilBERT, ale był od niego około ośmiokrotnie wolniejszy. Metody TF-IDF
    przetwarzały od 2,2 do 10,6 tys. wiadomości na sekundę, przy wyniku
    zewnętrznym niższym o 7,3--18,4 punktu procentowego. Przepustowość Qwen3,
    równa 7,3 wiadomości na sekundę, może ograniczać zastosowanie tego modelu
    jako jedynego filtra całego ruchu. Pomiary wykonane na jednym komputerze
    nie wystarczają jednak do wyceny konkretnego wdrożenia.
\end{enumerate}

\section{Ograniczenia pracy}

Eksperymenty przeprowadzono na publicznych korpusach, z których część powstała
wiele lat temu. Wiadomości z poszczególnych źródeł różnią się formatowaniem,
słownictwem i sposobem nadawania etykiet. Deduplikacja ograniczyła liczbę
powtórzeń, ale nie usunęła wszystkich cech pozwalających rozpoznać źródło
wiadomości. Wyników nie potwierdzono na bieżącym strumieniu poczty.

Trzy zbiory zewnętrzne nie były używane do treningu Qwen3, ale ich wyniki
służyły do wyboru konfiguracji i punktu kontrolnego. W metodach bazowych
parametry dobierano wyłącznie na części walidacyjnej, a wymienione zbiory
wykorzystano dopiero do oceny wybranego wariantu.

Każdy trening wykonano z jednym ziarnem losowym. Nie obliczano przedziałów
ufności i nie powtarzano tych samych konfiguracji, dlatego różnice rzędu
0,1--0,3 punktu procentowego mogą wynikać częściowo z losowości treningu lub
próbkowania. Dotyczy to zwłaszcza kolejności najlepszych sposobów dostrajania.

Zakres strojenia nie był jednakowy. Szesnaście konfiguracji i pośrednie punkty
kontrolne oceniono dla metody następnego tokenu. Pozostałe metody korzystały z
wybranej części tej przestrzeni, a dla odpowiedzi ustrukturyzowanej ukończono
dwie konfiguracje. Dla metod bazowych sprawdzono ograniczone zbiory parametrów
i jeden podział walidacyjny. Naive Bayes nie podlegał strojeniu, a DistilBERT
wytrenowano tylko z jednym zestawem parametrów.
Przedstawione wyniki dotyczą więc wyłącznie sprawdzonych konfiguracji.

Badanie objęło jeden model bazowy i jeden jego rozmiar. Pomiar wydajności
wykonano na jednym komputerze i dla jednej próbki. Nie sprawdzono zachowania
przy wielu równoczesnych żądaniach, opóźnień krańcowych, zużycia energii ani
kosztu pracy serwera. Zmierzone wartości pokazują względną różnicę między
metodami w badanym środowisku, ale nie opisują infrastruktury produkcyjnej.

Klasyfikacja binarna łączy spam reklamowy, phishing i wiadomości oszukańcze w
jedną klasę. Model nie określa rodzaju zagrożenia ani nie wyjaśnia, który
fragment wiadomości wpłynął na decyzję. Do modelu przekazywano tylko temat i
treść, bez reputacji nadawcy, pełnych nagłówków i mechanizmów uwierzytelniania
domeny.

\section{Możliwość wdrożenia}

Koszt inferencji porównano na wspólnej próbce 1000 wiadomości. Qwen3 z
adapterem LoRA przetwarzał 7,3 wiadomości na sekundę, a przetworzenie całej
próbki trwało 136,7 sekundy. Proces wykorzystał szczytowo 3466~MiB pamięci
operacyjnej, natomiast model bazowy i adapter zajmowały razem około 1526~MiB na
dysku. DistilBERT przetwarzał 57,9 wiadomości na sekundę, wykorzystywał
szczytowo 563~MiB pamięci i zajmował 256~MiB na dysku. Dla TF-IDF z Naive
Bayes przepustowość wyniosła około 10,6 tys. wiadomości na sekundę, szczytowe
użycie pamięci 337~MiB, a rozmiar modelu z wektoryzatorem 13~MiB.

Zmierzona przepustowość Qwen3 ogranicza możliwość wykorzystania go jako
jedynego filtra pierwszej linii bez skalowania infrastruktury. Szybszy sprzęt,
kwantyzacja i przetwarzanie wiadomości w partiach mogłyby poprawić
przepustowość. Na podstawie wykonanych
testów nie można oszacować kosztu konkretnego wdrożenia serwerowego.

DistilBERT osiągnął najlepszy wynik wśród metod bazowych i działał szybciej od
Qwen3, ale nadal był wielokrotnie wolniejszy od TF-IDF.

Koszt inferencji można byłoby ograniczyć przez zastosowanie układu kaskadowego.
Uwierzytelnianie nadawcy, reguły, reputacja i tani klasyfikator tekstu
obsługiwałyby przypadki jednoznaczne, a model językowy analizowałby tylko
pozostałe wiadomości. W pracy nie sprawdzono, czy Qwen3 poprawia wyniki właśnie
dla przykładów, co do których prostszy klasyfikator jest niepewny.

Model Qwen3-0.6B z adapterem LoRA jest prototypem klasyfikatora treści, a nie
kompletnym filtrem pocztowym. W systemie produkcyjnym jego wynik mógłby być
jednym z sygnałów używanych do podjęcia decyzji. Samodzielne odrzucanie
wiadomości wymagałoby dokładniejszej kontroli fałszywych alarmów i testów na
danych produkcyjnych.

\section{Kierunki dalszych badań}

Wyniki odpowiedzi ustrukturyzowanej trzeba uzupełnić o dwie brakujące
konfiguracje. Szersze przeszukanie parametrów Naive Bayes i DistilBERT
pozwoliłoby sprawdzić, czy ich obecne wyniki są bliskie najlepszym dla przyjętych
danych. Osobno trzeba dobrać progi klasyfikacji do założonego kosztu
\emph{false positive} i \emph{false negative}.

Kolejny eksperyment powinien rozdzielić wybór konfiguracji Qwen3 od testu na
zbiorach zewnętrznych. Można również przeprowadzić walidację
\emph{leave-one-source-out}, kolejno wyłączając z treningu każde źródło. Innym
wariantem jest podział czasowy, w którym do testu trafiają nowsze wiadomości.
Pomiar na niezbalansowanym strumieniu pozwoliłby
ocenić liczbę fałszywych alarmów przy proporcjach klas zbliżonych do ruchu
produkcyjnego. Osobne etykiety
spamu reklamowego, phishingu i oszustw pozwoliłyby sprawdzić, czy model
rozpoznaje także rodzaj zagrożenia.

Ocena wdrożeniowa wymaga porównania wariantów kwantyzacji, silników inferencji
i ustawień przetwarzania partiami na sprzęcie serwerowym. Oprócz średniej
przepustowości należy zmierzyć opóźnienie, użycie pamięci przy równoczesnych
żądaniach oraz koszt przetworzenia określonej liczby wiadomości. Osobnego testu
wymaga też układ kaskadowy z tanim klasyfikatorem pierwszego etapu.

W badanych warunkach dostrojony Qwen3-0.6B uzyskał wyższy wynik zbiorczy niż
metody bazowe, lecz wymagał znacznie więcej czasu i pamięci. Qwen3 mógłby
pełnić rolę klasyfikatora drugiego etapu, ale takiego układu nie oceniono w
eksperymentach.

 %%%%%%%%%%%%%%%%%%%%%%%%%%%%%%%%%%%%%
 %%%%%%%%%%%%%%%%%%%%%%%%%%%%%%%%%%%%%
\appendix % Dodatek
 %%%%%%%%%%%%%%%%%%%%%%%%%%%%%%%%%%%%%
 % Rozdziały dodatku
\include{tex/example}
 %%%%%%%%%%%%%%%%%%%%%%%%%%%%%%%%%%%%%
 %%%%%%%%%%%%%%%%%%%%%%%%%%%%%%%%%%%%%
\backmatter % Część końcowa
 %%%%%%%%%%%%%%%%%%%%%%%%%%%%%%%%%%%%%
 % Rozdziały części końcowej
\include{tex/notes}
 %%%%%%%%%%%%%%%%%%%%%%%%%%%%%%%%%%%%%
 %%%%%%%%%%%%%%%%%%%%%%%%%%%%%%%%%%%%%
 % UWAGA: Jeżeli będziesz tworzyć spis literatury ręcznie, tj. za pomocą otoczenia 
 %  'thebibliography', to pamiętaj — powinien zawierać tylko te publikacje, 
 %  do których odwołujesz się w pracy. 

\printbibliography  % Wyświetl spis literatury
 %%%%%%%%%%%%%%%%%%%%%%%%%%%%%%%%%%%%%
 %%%%%%%%%%%%%%%%%%%%%%%%%%%%%%%%%%%%%
\end{document}
